\ifdefined\gkver\else\def\gkver{student}\fi
\documentclass[\gkver]{gaokaozhenti}
\begin{document}

\chapter{1991年湖南云南海南}

\begin{enumerate}
\item
%% number: 1
%% paperName: 1991年湖南云南海南高考第 1 题 3 分
%% typeId: 1
%% type: 单选题
%% chapter: 运动和力
%% point: 动量定理
%% method: 无
%% score: 3
%% degree: 850
%% duplicateId: 0
%% body:
一个运动物体，从某时刻起仅受一给定的恒定阻力作用而逐渐减速，直到停止。这段运动时间由下列的哪个物理量完全决定\xzanswer{C}
\fourchoices[answer=C]
{物体的初速度}
{物体的初动能}
{物体的初动量}
{物体的质量}

%% answer: C
%% memo:
\memoanswer{
由动量定理得 $-Ft = 0 - mv_{0}$，解得 $t = \frac{mv_{0}}{F} = \frac{p_{0}}{F}$，可知运动的时间由物体的初动量决定，故 C 正确。
}

\item
%% number: 2
%% paperName: 1991年湖南云南海南高考第 2 题 3 分
%% typeId: 1
%% type: 单选题
%% chapter: 运动和力
%% point: 动量守恒定律
%% method: 无
%% score: 3
%% degree: 750
%% duplicateId: 0
%% body:
设 a、b 两小球相撞，碰撞前后都在同一直线上运动。若测得它们相撞前的速度为 $v_{a}$、$v_{b}$，相撞后的速度为 $v_{a}'$、$v_{b}'$，可知两球的质量之比 $\frac{m_{a}}{m_{b}}$ 等于\xzanswer{A}
\fourchoices[answer=A]
{$\frac{v_{b}'-v_{b}}{v_{a}-v_{a}'}$}
{$\frac{v_{a}'-v_{a}}{v_{b}'-v_{b}}$}
{$\frac{v_{a}'-v_{b}'}{v_{a}-v_{b}}$}
{$\frac{v_{a}-v_{a}'}{v_{b}'-v_{b}}$}

%% answer: A
%% memo:
\memoanswer{
两球碰撞过程时系统动量守恒，以 a 的初速度方向为正方向，由动量守恒定律得 $m_{a}v_{a}+m_{b}v_{b}=m_{a}v_{a}'+m_{b}v_{b}'$，解得 $\frac{m_{a}}{m_{b}}=\frac{v_{b}'-v_{b}}{v_{a}-v_{a}'}$，故 A 正确。
}

\item
%% number: 3
%% paperName: 1991年湖南云南海南高考第 3 题 3 分
%% typeId: 1
%% type: 单选题
%% chapter: 运动和力
%% point: 牛顿定律的应用
%% method: 无
%% score: 3
%% degree: 550
%% duplicateId: 0
%% body:
有三个光滑斜轨道 1、2、3，它们的倾角依次是 $60^{\circ}$、$45^{\circ}$ 和 $30^{\circ}$，这些轨道交于 O 点。现有位于同一竖直线上的 3 个小物体甲、乙、丙，分别沿这 3 个轨道同时从静止自由下滑，如图所示。物体滑到 O 点的先后顺序是\xzanswer{B}
\onepicture[w=4cm]{figs/03.png}
\fourchoices[answer=B]
{甲最先，乙稍后，丙最后}
{乙最先，然后甲和丙同时到达}
{甲、乙、丙同时到达}
{乙最先，甲稍后，丙最后}

%% answer: B
%% memo:
\memoanswer{
设斜轨道底边的长度为 $l$，斜面的倾角为 $\alpha$，则斜轨道的长度为 $x=\frac{l}{\cos\alpha}$，由牛顿第二定律知物体下滑的加速度为 $a=\frac{mg\sin\alpha}{m}=g\sin\alpha$，又 $x=\frac{1}{2}at^{2}$，代入数据得 $\frac{l}{\cos\alpha}=\frac{1}{2}g\sin\alpha t^{2}$，得 $t=\sqrt{\frac{2l}{g\sin\alpha\cos\alpha}}=\sqrt{\frac{4l}{g\sin2\alpha}}$，由数学知识得知，$\sin(2\times60^{\circ})=\sin(2\times30^{\circ})$，则甲和丙运动的时间相等，同时达到斜轨道的底端 O 点，又 $\sin(2\times45^{\circ})=1$ 最大，则乙运动时间最短，乙最先到达 O 点，故 B 正确。
}

\item
%% number: 4
%% paperName: 1991年湖南云南海南高考第 4 题 3 分
%% typeId: 1
%% type: 单选题
%% chapter: 气体性质
%% point: 液柱移动定性分析
%% method: 无
%% score: 3
%% degree: 650
%% duplicateId: 0
%% body:
两端封闭的均匀玻璃管，水平放置，管内有一小段水银将气体分成左右两部分，体积为 $V_{\text{左}}$ 和 $V_{\text{右}}$，它们的温度均为 $T_{1}$，现将两边气体的温度同时缓慢地升高到 $T_{2}$，在升温过程中\xzanswer{D}
\fourchoices[answer=D]
{若 $V_{\text{左}} > V_{\text{右}}$，则水银柱将向左移动}
{若 $V_{\text{左}} > V_{\text{右}}$，则水银柱将向右移动}
{只有当 $V_{\text{左}} = V_{\text{右}}$ 时，水银柱才能保持不动}
{无论 $V_{\text{左}}$、$V_{\text{右}}$ 大小如何，水银柱都保持不动}

%% answer: D
%% memo:
\memoanswer{
假设水银柱不动，即 $V$ 不变，有 $\frac{p}{T}=\frac{p+\Delta p}{T+\Delta T}$，得 $\Delta p=\frac{p\Delta T}{T}$，一开始两封闭气体的压强相等，所以升温后，两封闭气体的压强增加量相等，因此水银柱都保持不动，故 D 正确。
}

\item
%% number: 5
%% paperName: 1991年湖南云南海南高考第 5 题 3 分
%% typeId: 1
%% type: 单选题
%% chapter: 电场
%% point: 静电
%% method: 无
%% score: 3
%% degree: 750
%% duplicateId: 0
%% body:
绝缘细线上端固定，下端悬挂一个轻质小球 a，a 的表面镀有铝膜。在 a 近旁有一绝缘金属球 b。开始时 a、b 都不带电，如图所示。现使 b 带电，则\xzanswer{B}
\onepicture[w=3.5cm]{figs/05.png}
\fourchoices[answer=B]
{b 将吸引 a，吸住后不放开}
{b 先吸引 a，接触后又把 a 排斥开}
{a、b 之间不发生相互作用}
{b 立即把 a 排斥开}

%% answer: B
%% memo:
\memoanswer{
带电物体能够吸引轻小物体，a 与 b 接触后，a 球与 b 球带上了同种电荷，相互排斥而分开，故 B 正确。
}

\item
%% number: 6
%% paperName: 1991年湖南云南海南高考第 6 题 3 分
%% typeId: 1
%% type: 单选题
%% chapter: 电场
%% point: 带电粒子的偏转
%% method: 无
%% score: 3
%% degree: 650
%% duplicateId: 0
%% body:
一个带负电的小球，受水平方向的匀强电场力和重力的作用，由静止开始运动。不计空气阻力，设坐标轴如下图，$x$ 轴的正方向与电场方向一致，$y$ 轴向下，原点在小球起始位置。在下列四个图示中，哪个图可能表示此小球的运动轨迹？\xzanswer{D}
\fourchoices[answer=D, ispicture=true, h=2.4cm]
{figs/06a.png}
{figs/06b.png}
{figs/06c.png}
{figs/06d.png}

%% answer: D
%% memo:
\memoanswer{
由于小球初速度为 0，且受到恒力的作用，因此小球做直线运动，其中电场力与重力合成的等效重力方向指向左下角，所以小球将沿等效重力的方向做匀加速直线运动，故 D 正确。
}

\item
%% number: 7
%% paperName: 1991年湖南云南海南高考第 7 题 3 分
%% typeId: 1
%% type: 单选题
%% chapter: 电路
%% point: 动态分析
%% method: 无
%% score: 3
%% degree: 650
%% duplicateId: 0
%% body:
如图所示的 4 个电路中，a、b 两点与一稳压直流电源相接。当滑动变阻器的滑动点 O 向 d 点移动一段距离时，哪一个电路中的电流表读数会变小\xzanswer{B}
\fourchoices[answer=B, ispicture=true, h=2.2cm]
{figs/07a.png}
{figs/07b.png}
{figs/07c.png}
{figs/07d.png}

%% answer: B
%% memo:
\memoanswer{
电路中电流表的示数要变小，由 $I=\frac{U}{R}$ 可知，在 a、b 两点间电压不变的条件下，电路中的阻值应变大；A、C 选项中滑动变阻器接入的电阻是 Od 段，滑片 O 向 d 端移动一段距离时，电路中的阻值减小，电流变大，不符合条件，故 A、C 错误；B 选项中滑动变阻器接入的电阻是 Oc 段，滑片 O 向 d 端移动一段距离时，电路中的阻值变大，电流变小，符合条件，故 B 正确；D 选项中滑动变阻器接入的电阻是 Od 段，滑片 O 向 d 端移动一段距离时，电路中的阻值减小，电流变大，不符合条件，故 D 错误。
}

\item
%% number: 8
%% paperName: 1991年湖南云南海南高考第 8 题 3 分
%% typeId: 1
%% type: 单选题
%% chapter: 光学
%% point: 透镜
%% method: 无
%% score: 3
%% degree: 750
%% duplicateId: 0
%% body:
在一薄透镜左侧有一光源 AB，在透镜右侧的光屏 S 上可得到光源的像。现移去光屏，从光源上 B 点发出的如图所示的光线 BP 经过透镜后\xzanswer{B}
\onepicture[w=7cm]{figs/08.png}
\fourchoices[answer=B]
{将与主光轴交于原光屏 S 的右侧}
{将与主光轴交于透镜与原光屏 S 之间}
{将平行于主光轴}
{其反向延长线将与主光轴交于透镜的左侧}

%% answer: B
%% memo:
\memoanswer{
在一透镜左侧有一光源 AB，在透镜右侧的光屏 S 上得到光源清晰的像，则像是倒立的实像，A 点在主光轴的上方，成像应在主光轴的下方，A 点发出的光经透镜折射后，射向主光轴的下方，与主轴交于透镜与原光屏 S 之间，故 B 正确。
}

\item
%% number: 9
%% paperName: 1991年湖南云南海南高考第 9 题 3 分
%% typeId: 1
%% type: 单选题
%% chapter: 光学
%% point: 光的干涉
%% method: 无
%% score: 3
%% degree: 750
%% duplicateId: 0
%% body:
马路上积水表面的油膜呈现彩色图样。这是由于\xzanswer{D}
\fourchoices[answer=D]
{光的色散}
{光的衍射}
{光的漫反射}
{光的干涉}

%% answer: D
%% memo:
\memoanswer{
题中叙述的是薄膜干涉的一个实例，其产生原因是薄膜（油膜）厚度不均匀性造成的。白光中的某一单色光照射到油膜上，从不均匀油膜前后表面反射回来的反射光，互相加强就出现了这种单色光的亮纹，在另一厚度的地方，另一波长的反射光相互加强，就形成另一单色光的亮纹，这样在油膜上呈现出彩色图像，这就是薄膜干涉现象，故 D 正确。
}

\item
%% number: 10
%% paperName: 1991年湖南云南海南高考第 10 题 3 分
%% typeId: 1
%% type: 单选题
%% chapter: 原子和原子核
%% point: 半衰期
%% method: 无
%% score: 3
%% degree: 750
%% duplicateId: 0
%% body:
某放射性同位素样品，在 21 天里衰减掉 $\frac{7}{8}$，它的半衰期是\xzanswer{C}
\fourchoices[answer=C]
{3 天}
{5.25 天}
{7 天}
{10.5 天}

%% answer: C
%% memo:
\memoanswer{
由半衰期的定义知 $m=m_{0}\left(\frac{1}{2}\right)^{\frac{t}{T}}$，即 $\frac{1}{8}=\left(\frac{1}{2}\right)^{\frac{t}{T}}$，可得 $\frac{t}{T}=3$，即 $T=\frac{t}{3}=7$ 天，故 C 正确。
}

\item
%% number: 11
%% paperName: 1991年湖南云南海南高考第 11 题 4 分
%% typeId: 2
%% type: 多选题
%% chapter: 原子和原子核
%% point: 人工核反应
%% method: 无
%% score: 4
%% degree: 750
%% duplicateId: 0
%% body:
在下列 4 个核反应式中，X 表示中子的是哪些\xzanswer{BCD}
\fourchoices[answer=BCD]
{\ce{^{14}_{7}N + ^{4}_{2}He -> ^{17}_{8}O + X}}
{\ce{^{27}_{13}Al + ^{4}_{2}He -> ^{30}_{15}P + X}}
{\ce{^{2}_{1}H + ^{3}_{1}H -> ^{4}_{2}He + X}}
{\ce{^{235}_{92}U + X -> ^{90}_{38}Sr + ^{136}_{54}Xe + 10^{1}_{0}n}}

%% answer: BCD
%% memo:
\memoanswer{
核反应方程满足核电荷数和质量数守恒，应用两守恒列等式，可计算得 A 项的 X 为 \ce{^{1}_{1}H}，B、C、D 项的 X 为 \ce{^{1}_{0}n}，故 B、C、D 正确。
}

\item
%% number: 12
%% paperName: 1991年湖南云南海南高考第 12 题 4 分
%% typeId: 1
%% type: 单选题
%% chapter: 电场
%% point: 电容
%% method: 无
%% score: 4
%% degree: 650
%% duplicateId: 0
%% body:
一平行板电容器的两个极板分别接在电池组的正极和负极上，若使两极板之间的距离增大，则\xzanswer{D}
\fourchoices[answer=D]
{两极板间匀强电场的电场强度保持不变}
{电容器所带的电量保持不变}
{电容器的带电量和电势差的比值保持不变}
{电容器的电容减少}

%% answer: D
%% memo:
\memoanswer{
由题意知，平行板电容器接在电源两端，故电容器的电压保持不变，即 $U$ 不变，又 $E=\frac{U}{d}$，当 $d$ 增大时，$E$ 减小，故 A 错误；由平行板电容器得决定式 $C=\frac{\varepsilon S}{4\pi kd}$ 知，又 $Q=CU$，当 $d$ 增大时，$C$ 减小，$Q$ 减小，故 B、C 错误，D 正确。
}

\item
%% number: 13
%% paperName: 1991年湖南云南海南高考第 13 题 4 分
%% typeId: 2
%% type: 多选题
%% chapter: 交流电
%% point: 交流电
%% method: 无
%% score: 4
%% degree: 650
%% duplicateId: 0
%% body:
在相同的时间内，某正弦交流电通过一阻值为 $100\UO$ 的电阻产生的热量，与一电流强度为 $3\UA$ 的直流电通过同一阻值的电阻产生的热量相等，则\xzanswer{AC}
\fourchoices[answer=AC]
{此交流电的电流强度的有效值为 $3\UA$，最大值为 $3\sqrt{2}\UA$}
{此交流电的电流强度的有效值为 $3\sqrt{2}\UA$，最大值为 $6\UA$}
{电阻两端的交流电电压的有效值为 $300\UV$，最大值为 $300\sqrt{2}\UV$}
{电阻两端的交流电电压的有效值为 $300\sqrt{2}\UV$，最大值为 $600\UV$}

%% answer: AC
%% memo:
\memoanswer{
在相同的时间内，某正弦交流电通过某一电阻产生的热量与一直流电通过同一阻值的电阻产生的热量相等时，此直流电的数值就等于交流电的有效值，所以此交流电电流的有效值为 $3\UA$，又正弦交流电电流的有效值 $I = \frac{I_{m}}{\sqrt{2}}$，所以交流电电流的最大值为 $3\sqrt{2}\UA$，故 A 正确；电阻两端的交流电电压的有效值 $U = IR = 300\UV$，交流电电压的最大值 $U_{m} = \sqrt{2}U = 300\sqrt{2}\UV$，故 C 正确。
}

\item
%% number: 14
%% paperName: 1991年湖南云南海南高考第 14 题 4 分
%% typeId: 2
%% type: 多选题
%% chapter: 电磁感应
%% point: 线圈过有界磁场
%% method: 无
%% score: 4
%% degree: 550
%% duplicateId: 0
%% body:
边长为 $h$ 的正方形金属导线框从图示的初始位置由静止开始下落，通过一匀强磁场区域。磁场方向是水平的，且垂直于线框平面。磁场区宽度等于 $H$，上下边界如图中水平虚线所示，$H > h$。从线框开始下落到完全穿过磁场区的整个过程中\xzanswer{CD}
\onepicture[w=4.5cm]{\includesvg{figs/14}}
\fourchoices[answer=CD]
{线框中总是有感应电流存在}
{线框受到的磁场力的合力的方向有时向上，有时向下}
{线框运动的方向始终是向下的}
{线框速度的大小不一定总是在增加}

%% answer: CD
%% memo:
\memoanswer{
线框在进入和穿出磁场的过程中有感应电流，线框全部处于磁场中时，磁通量不变，无感应电流，故 A 错误；线框进入磁场时磁通量增加，离开磁场时磁通量减小，由楞次定律可知，线框受到磁场力的合力方向总向上，故 B 错误；从线框开始下落到完全穿过磁场区域的全过程中，线框运动方向始终向下，故 C 正确；从下边开始离开磁场时，当安培力大于其重力时，线框会做减速运动，故 D 正确。
}

\item
%% number: 15
%% paperName: 1991年湖南云南海南高考第 15 题 4 分
%% typeId: 2
%% type: 多选题
%% chapter: 交流电
%% point: LC振荡回路
%% method: 无
%% score: 4
%% degree: 650
%% duplicateId: 0
%% body:
在 LC 回路产生电磁振荡的过程中\xzanswer{BC}
\fourchoices[answer=BC]
{电容器放电完毕的时刻，磁场能最小}
{回路中电流强度达到最大的时刻，磁场能最大}
{电容器极板上电荷最多的时刻，电场能最大}
{回路中电流强度最小的时刻，电场能最小}

%% answer: BC
%% memo:
\memoanswer{
电容器放电完毕时，带电荷量为零，电场能为零，电路中的感应电流最大，磁场能最大，故 A 错误、B 正确；电容器充电完毕时，带电荷量最大，电场能达到最大，而磁场能为零，回路中感应电流 $i = 0$，故 C 正确；回路中电流值最小时刻，磁场能最小，而电容器的带电荷量最多，电场能最大，故 D 错误。
}

\item
%% number: 16
%% paperName: 1991年湖南云南海南高考第 16 题 4 分
%% typeId: 2
%% type: 多选题
%% chapter: 气体性质
%% point: 分子力
%% method: 无
%% score: 4
%% degree: 650
%% duplicateId: 0
%% body:
用 $r$ 表示两个分子间的距离，$E_{p}$ 表示两个分子间相互作用的势能。当 $r = r_{0}$ 时两分子间斥力等于引力。设两分子相距很远时 $E_{p} = 0$\xzanswer{AB}
\fourchoices[answer=AB]
{当 $r > r_{0}$ 时，$E_{p}$ 随 $r$ 的增大而增加}
{当 $r < r_{0}$ 时，$E_{p}$ 随 $r$ 的减小而增加}
{当 $r > r_{0}$ 时，$E_{p}$ 不随 $r$ 而变}
{当 $r = r_{0}$ 时，$E_{p} = 0$}

%% answer: AB
%% memo:
\memoanswer{
当 $r > r_{0}$ 时，分子力表现为引力，当 $r$ 的增加时，分子力做负功，则 $E_{p}$ 增加，故 A 正确、C 错误；当 $r < r_{0}$ 时，分子力表现为斥力，当 $r$ 的减小时，分子力做负功，则 $E_{p}$ 增加，故 B 正确；当 $r = r_{0}$ 时，$E_{p}$ 最小，而不是零，故 D 错误。
}

\item
%% number: 17
%% paperName: 1991年湖南云南海南高考第 17 题 4 分
%% typeId: 2
%% type: 多选题
%% chapter: 光学
%% point: 透镜
%% method: 无
%% score: 4
%% degree: 550
%% duplicateId: 0
%% body:
一物体经过一个薄透镜成像，以下 4 条论断中，哪几条是完全正确的\xzanswer{CD}
\fourchoices[answer=CD]
{经凸透镜所成实像总是倒立的，放大的，且与物体分别位于透镜的两侧}
{经凹透镜只能得到虚像，此像是放大的还是缩小的，取决于物体位置是在焦点以内还是在焦点以外}
{经凸透镜得到的虚像总是正立的和放大的}
{当物体到凸透镜的距离小于焦距时，在透镜另一侧任何位置的屏上都得不到物体的像}

%% answer: CD
%% memo:
\memoanswer{
经凸透镜所成的实像总是倒立的，且为缩小、等大或放大的实像，故 A 错误；当物体为虚物时，凹透镜到虚物的距离为一倍焦距以内时，成正立、放大的实像，像与物在透镜的同侧，故 B 错误；由凸透镜成像规律知，凸透镜得到的虚像都是正立且放大的，故 C 正确；当物体到凸透镜的距离小于焦距时，成正立、放大的虚像，故屏上得不到物体的像，故 D 正确。
}

\item
%% number: 18
%% paperName: 1991年湖南云南海南高考第 18 题 4 分
%% typeId: 2
%% type: 多选题
%% chapter: 光学
%% point: 光的电磁说
%% method: 无
%% score: 4
%% degree: 650
%% duplicateId: 0
%% body:
某种无色透明玻璃对于真空中波长为 $0.60\Uum$ 的单色光的折射率是 1.50。则\xzanswer{AB}
\fourchoices[answer=AB]
{这种光的频率是 $5.0\times10^{14}\UHz$}
{这种光在该玻璃中的传播速度是 $2.0\times10^{8}\Ums$}
{这种光在该玻璃中的波长是 $0.90\Uum$}
{对于真空中波长为 $0.40\Uum$ 的单色光，该玻璃的折射率略小于 1.50}

%% answer: AB
%% memo:
\memoanswer{
由 $c = \lambda f$ 得光的频率为 $f = \frac{c}{\lambda} = \frac{3\times10^{8}\Ums}{0.60\times10^{-6}\Um} = 5.0\times10^{14}\UHz$，故 A 正确；光在玻璃中的频率等于在真空中的频率，光在玻璃中的传播速度为 $v = \frac{c}{n} = \frac{3\times10^{8}\Ums}{1.5} = 2\times10^{8}\Ums$，由 $v = \lambda f$ 得光在该玻璃中的波长 $\lambda = \frac{v}{f} = \frac{2\times10^{8}\Ums}{5\times10^{14}\Us^{-1}} = 0.4\times10^{-6}\Um = 0.4\Uum$，故 B 正确、C 错误；光的波长越小，频率越大，玻璃对光的折射率越大，所以对于真空中波长为 $0.4\Uum$ 的单色光，该玻璃的折射率大于 1.50，故 D 错误。
}

\item
%% number: 19
%% paperName: 1991年湖南云南海南高考第 19 题 4 分
%% typeId: 2
%% type: 多选题
%% chapter: 运动和力
%% point: 连接体
%% method: 无
%% score: 4
%% degree: 550
%% duplicateId: 0
%% body:
如图所示，水平地面上有两个完全相同的木块 A、B，在水平推力 $F$ 作用下运动。用 $F_{AB}$ 代表 A、B 间的相互作用力\xzanswer{BD}
\onepicture[w=5cm]{\includesvg{figs/19}}
\fourchoices[answer=BD]
{若地面是完全光滑的，则 $F_{AB} = F$}
{若地面是完全光滑的，则 $F_{AB} = \frac{1}{2}F$}
{若地面的摩擦系数为 $\mu$，则 $F_{AB} = F$}
{若地面的摩擦系数为 $\mu$，则 $F_{AB} = \frac{1}{2}F$}

%% answer: BD
%% memo:
\memoanswer{
设两木块的质量均为 $m$，若地面是完全光滑的，对整体，由牛顿第二定律得 $F = 2ma$；对 B，由牛顿第二定律得 $F_{AB} = ma = \frac{F}{2}$，故 A 错误、B 正确；若地面动摩因数为 $\mu$，对整体，由牛顿第二定律得 $F - 2\mu mg = 2ma$，对 B，由牛顿第二定律得 $F_{AB} - \mu mg = ma$，解得 $F_{AB} = \frac{F}{2}$，故 C 错误、D 正确。
}

\item
%% number: 20
%% paperName: 1991年湖南云南海南高考第 20 题 4 分
%% typeId: 1
%% type: 单选题
%% chapter: 光学
%% point: 全反射
%% method: 无
%% score: 4
%% degree: 550
%% duplicateId: 0
%% body:
abc 为一全反射棱镜，它的主截面是等腰直角三角形，如图所示。一束白光垂直入射到 ac 面上，在 ab 面上发生全反射。若光线入射点 O 的位置保持不变，改变光线的入射方向（不考虑自 bc 面反射的光线）\xzanswer{A}
\onepicture[w=5cm]{figs/20.png}
\fourchoices[answer=A]
{使入射光按图中所示的顺时针方向逐渐偏转，如果有色光射出 ab 面，则红光将首先射出}
{使入射光按图中所示的顺时针方向逐渐偏转，如果有色光射出 ab 面，则紫光将首先射出}
{使入射光按图中所示的逆时针方向逐渐偏转，红光将首先射出 ab 面}
{使入射光按图中所示的逆时针方向逐渐偏转，紫光将首先射出 ab 面}

%% answer: A
%% memo:
\memoanswer{
白光由从红到紫七种色光组成，同一种介质对它们的折射率，从红光到紫光逐渐增大。在同一种介质中产生全反射，它们的临界角不同，由 $\sin C = \frac{1}{n}$ 知，$n$ 越小，$C$ 越大。红光折射率最小，则临界角最大。光垂直入射到 ac 面，在 ab 面发生全反射，则临界角 $C \leqslant 45^{\circ}$。当光沿顺时针方向偏转入射，其入射角 $C$ 减小，如图 1 所示，首先减小到红光临界角以下，红光先射出 ab 面，故 A 正确、B 错误；当光沿逆时针方向偏转入射，如图 2 所示，其入射角增大，不可能有光线在 ab 面上射出，故 C、D 错误。

\onepicture[w=5.5cm]{figs/20_an1.png}

\onepicture[w=5.5cm]{figs/20_an2.png}
}

\item
%% number: 21
%% paperName: 1991年湖南云南海南高考第 21 题 4 分
%% typeId: 3
%% type: 填空题
%% chapter: 功和能
%% point: 交通工具启动
%% method: 无
%% score: 4
%% degree: 850
%% duplicateId: 0
%% body:
汽车在水平的公路上沿直线匀速行驶，当速度为 $18\Ums$ 时，其输出功率为 $72\UkW$，汽车所受到的阻力是 \tkanswer[2.5cm]{$4.0\times10^{3}$}$\UN$。
%% answer: 
\jdanswer{
$4.0\times10^{3}\UN$
}

%% memo:
\memoanswer{
由题意知，汽车匀速运动的速度为 $v = 18\Ums$，汽车的输出功率为 $P = 7.2 \times 10^4\UW$，由 $P = Fv$ 知汽车匀速行驶的阻力 $f = F = \frac{P}{v} = \frac{7.2 \times 10^4}{18} = 4000\UN$。
}

\item
%% number: 22
%% paperName: 1991年湖南云南海南高考第 22 题 4 分
%% typeId: 3
%% type: 填空题
%% chapter: 力物体的平衡
%% point: 三力平衡
%% method: 无
%% score: 4
%% degree: 750
%% duplicateId: 0
%% body:
一物块在与水平面成 $\theta$ 角的斜面上匀速下滑。物块与斜面间的滑动摩擦系数是 \tkanswer[1.5cm]{$\tan\theta$}。
%% answer: 
\jdanswer{
$\tan\theta$
}

%% memo:
\memoanswer{
因物体沿斜面匀速下滑，物块受力平衡，有 $mg\sin\theta = \mu mg\cos\theta$，解得 $\mu = \tan\theta$。
}

\item
%% number: 23
%% paperName: 1991年湖南云南海南高考第 23 题 4 分
%% typeId: 3
%% type: 填空题
%% chapter: 运动和力
%% point: 动量和冲量
%% method: 无
%% score: 4
%% degree: 750
%% duplicateId: 0
%% body:
设质量为 $m$ 的小球从离地面高为 $h$ 处自由下落，着地后又被弹回原处。在小球与地面碰撞的过程中，小球所受的冲量的大小为 \tkanswer[2.5cm]{$2m\sqrt{2gh}$}。
%% answer: 
\jdanswer{
$2m\sqrt{2gh}$
}

%% memo:
\memoanswer{
取竖直向下方向为正方向，小球落地时的速度为 $v = \sqrt{2gh}$，着地后被弹回时的速度为 $v' = -v = -\sqrt{2gh}$，由动量定理得，地面对小球的冲量 $I = mv' - mv = -2mv = -2m\sqrt{2gh}$，故方向竖直向上，大小为 $2m\sqrt{2gh}$。
}

\item
%% number: 24
%% paperName: 1991年湖南云南海南高考第 24 题 4 分
%% typeId: 3
%% type: 填空题
%% chapter: 曲线运动
%% point: 天体运动
%% method: 无
%% score: 4
%% degree: 650
%% duplicateId: 0
%% body:
某行星的一颗小卫星在半径为 $r$ 的圆轨道上绕该行星运行，运行的周期是 $T$，已知万有引力恒量 $G$，这个行星的质量 $M =$ \tkanswer[2.5cm]{$\frac{4\pi^{2}r^{3}}{GT^{2}}$}。
%% answer: 
\jdanswer{
$\frac{4\pi^{2}r^{3}}{GT^{2}}$
}

%% memo:
\memoanswer{
由万有引力提供向心力得 $\frac{GMm}{r^{2}}=m\frac{4\pi^{2}}{T^{2}}r$，则行星的质量 $M=\frac{4\pi^{2}r^{3}}{GT^{2}}$。
}

\item
%% number: 25
%% paperName: 1991年湖南云南海南高考第 25 题 4 分
%% typeId: 3
%% type: 填空题
%% chapter: 电场
%% point: 带电粒子的圆周运动
%% method: 无
%% score: 4
%% degree: 650
%% duplicateId: 0
%% body:
在场强为 $E$ 的水平方向的匀强电场中，有一质量不计的轻杆，可绕杆的一端点 O 自由转动，另一端连一质量为 $m$ 的带正电的小球，把杆拉成水平后由静止释放（如图）。若小球达到最低位置时速度恰好为零，则小球所带的电量是 \tkanswer[2.5cm]{$\frac{mg}{E}$}。
\onepicture[align=right, w=4.5cm]{figs/25.png}
%% answer: 
\jdanswer{
$\frac{mg}{E}$
}

%% memo:
\memoanswer{
设杆长为 $L$，对小球分析可知，小球受重力和电场力的作用，由动能定理可得 $mgL - qEL = 0$，解得 $q = \frac{mg}{E}$。
}

\item
%% number: 26
%% paperName: 1991年湖南云南海南高考第 26 题 4 分
%% typeId: 3
%% type: 填空题
%% chapter: 运动和力
%% point: 瞬时加速度
%% method: 无
%% score: 4
%% degree: 650
%% duplicateId: 0
%% body:
一条轻弹簧和一根细线共同拉住一个质量为 $m$ 的小球，平衡时细线是水平的，弹簧与竖直方向的夹角是 $\theta$，如图所示。若突然剪断细线，则在刚剪断的瞬时，弹簧拉力的大小是 \tkanswer[2.5cm]{$\frac{mg}{\cos\theta}$}，小球加速度的方向与竖直方向的夹角等于 \tkanswer[2cm]{$90^{\circ}$}。
\onepicture[align=right, w=4cm]{figs/26.png}
%% answer: 
\jdanswer{
$\frac{mg}{\cos\theta}$，$90^{\circ}$
}

%% memo:
\memoanswer{
剪断细线的瞬间，细线的拉力消失，而弹簧还是保持原来的长度，弹力与剪断前相同，此时，小球受到重力和弹簧弹力的合力方向水平向右，与剪断细线前细线的拉力大小相等、方向相反，如图所示，即 $F_{\text{弹}} = \frac{mg}{\cos\theta}$，$F_{\text{合}} = mg\tan\theta$，由牛顿第二定律得，小球的加速度 $a = \frac{F_{\text{合}}}{m} = g\tan\theta$，方向水平向右。

\onepicture[w=3cm]{figs/26_an.png}
}

\item
%% number: 27
%% paperName: 1991年湖南云南海南高考第 27 题 4 分
%% typeId: 3
%% type: 填空题
%% chapter: 直线运动
%% point: 螺旋测微器
%% method: 无
%% score: 4
%% degree: 850
%% duplicateId: 0
%% body:
用螺旋测微器测量一圆柱形工件的直径时的示数如图，此示数应读为 \tkanswer[2cm]{$6.926$}$\Umm$。
\onepicture[align=right, w=4.5cm]{figs/27.png}
%% answer: 
\jdanswer{
$6.926\Umm$（$6.924 \sim 6.928\Umm$）
}

%% memo:
\memoanswer{
固定读数为 $6.5\Umm$，可动读数为 $42.6 \times 0.01\Umm = 0.426\Umm$，故读数为 $6.5\Umm + 0.426\Umm = 6.926\Umm$。
}

\item
%% number: 28
%% paperName: 1991年湖南云南海南高考第 28 题 4 分
%% typeId: 3
%% type: 填空题
%% chapter: 电路
%% point: 多用表的使用
%% method: 无
%% score: 4
%% degree: 650
%% duplicateId: 0
%% body:
一万用表的欧姆挡有 4 挡，分别为 $\times1\UO$，$\times10\UO$，$\times100\UO$，$\times1000\UO$。现用它来测一未知电阻值。当用 $\times100\UO$ 挡测量时，发现指针的偏转角度很小。为了使测量结果更准确，测量前应进行如下两项操作，先 \tkanswer[3.5cm]{把选择开关拨到 $\times1000\UO$ 挡}，接着 \tkanswer[2cm]{欧姆调零}，然后再测量并读数。
%% answer: 
\jdanswer{
把选择开关拨到 $\times1000\UO$ 挡；欧姆调零
}

%% memo:
\memoanswer{
当用 $\times100\UO$ 挡测量时，发现指针的偏转角度很小，说明待测电阻很大，故选用 $\times1000\UO$ 挡，使用其他欧姆挡测量时需要重新欧姆调零，测量的数值才准确。
}

\item
%% number: 29
%% paperName: 1991年湖南云南海南高考第 29 题 4 分
%% typeId: 3
%% type: 填空题
%% chapter: 力物体的平衡
%% point: 力矩平衡
%% method: 无
%% score: 4
%% degree: 450
%% duplicateId: 0
%% body:
一根长为 $L$ 的木头，质量大于 100 kg，其重心 O 在离大头为 $\frac{L}{3}$ 的地方。甲、乙二人同时各扛一端将木头扛起。此后丙又在木头全长的中点 N 处向上扛，用力 $f_{\text{丙}} = 300\UN$，如图所示。由于丙的参加，甲的负重减轻了 \tkanswer[2cm]{$150$}$\UN$，乙的负重减轻了 \tkanswer[2cm]{$150$}$\UN$。
\onepicture[align=right, w=5cm]{figs/29.png}
%% answer: 
\jdanswer{
$150\UN$；$150\UN$
}

%% memo:
\memoanswer{
丙参加作用前，设木头的质量为 $100\Ukg$，以右端为支点，由力矩平衡得

$f_{\text{甲}} \times L = 100\Ukg \times 10\UN/\Ukg \times \frac{2}{3}L$，解得 $f_{\text{甲}} = \frac{2000}{3}\UN$，

丙参加作用前，以左端为支点，由力矩平衡得

$f_{\text{乙}} \times L = 100\Ukg \times 10\UN/\Ukg \times \frac{1}{3}L$，解得 $f_{\text{乙}} = \frac{1000}{3}\UN$，

丙参加作用后，以右端为支点，由力矩平衡得

$f_{\text{甲}}'L + 300\UN \times \frac{1}{2}L = 100\Ukg \times 10\UN/\Ukg \times \frac{2}{3}L$，

解得 $f_{\text{甲}}' = \frac{1550}{3}\UN$，

丙参加作用后，以左端为支点，由力矩平衡得

$f_{\text{乙}}'L + 300\UN \times \frac{1}{2}L = 100\Ukg \times 10\UN/\Ukg \times \frac{1}{3}L$，

解得 $f_{\text{乙}}' = \frac{550}{3}\UN$，

故 $\Delta f_{\text{甲}} = f_{\text{甲}} - f_{\text{甲}}' = 150\UN$，$\Delta f_{\text{乙}} = f_{\text{乙}} - f_{\text{乙}}' = 150\UN$。
}

\item
%% number: 30
%% paperName: 1991年湖南云南海南高考第 30 题 4 分
%% typeId: 3
%% type: 填空题
%% chapter: 机械波
%% point: 波与振动图象
%% method: 无
%% score: 4
%% degree: 450
%% duplicateId: 0
%% body:
一简谐波沿 $x$ 轴正方向传播。已知轴上 $x_{1} = 0$ 和 $x_{2} = 1\Um$ 两处的振动图线分别如图 1 和图 2 所示。又知此波的波长大于 $1\Um$，则此波的传播速度 $v =$ \tkanswer[2.5cm]{$333$}$\Ums$。
\onepicture[align=right, w=8cm]{figs/30.png}
%% answer: 
\jdanswer{
$333\Ums$
}

%% memo:
\memoanswer{
由题意知，振幅为 $5 \times 10^{-3}\Um$，周期为 $4 \times 10^{-3}\Us$，$t = 0$ 时，$x_{1} = 0$ 的质点位于位移最大处，$x_{2} = 1\Um$ 的质点位于平衡位置，振动方向竖直向下。又波长大于 $1\Um$，两质点的间距小于一个波长，故两点的间距为 $\frac{3}{4}\lambda$，故简谐波的波长为 $\frac{4}{3}\Um$，由波速公式得 $v = \frac{\lambda}{T} = \frac{\frac{4}{3}\Um}{4 \times 10^{-3}\Us} = \frac{1000}{3}\Ums \approx 333\Ums$。
}

\item
%% number: 31
%% paperName: 1991年湖南云南海南高考第 31 题 6 分
%% typeId: 4
%% type: 实验题
%% chapter: 电路
%% point: 电路实验
%% method: 无
%% score: 6
%% degree: 550
%% duplicateId: 0
%% body:
用实验原理图（a）所示的电路测量电流计 G 的内阻。G 的量程是 $300\UmuA$，内阻约 $100\UO$。
\begin{enumerate}
	\item
	图（b）中已画出进行实验所需的器材，请按原理图要求连接成实验电路。
	\item
	实验室中有如下不同规格的电源与变阻器：

	A．电源（电动势 $2\UV$）

	B．电源（电动势 $4\UV$）

	C．滑动变阻器（阻值范围 $0\sim100\UO$）

	D．滑动变阻器（阻值范围 $0\sim10\UkO$）

	为了保证上述实验能顺利进行，电源应选用 \tkanswer[1.2cm]{B}（填字母代号），滑动变阻器应选用 \tkanswer[1.2cm]{D}（填字母代号）。
\end{enumerate}
\onepicture[align=right, w=8cm]{figs/31.png}
%% answer: （1）实物图连接如图所示；（2）B；D
\jdanswer{
\begin{enumerate}
	\item
	实物图连接如图所示。

	\onepicture[w=7cm]{figs/31_an.png}

	\item
	B；D
\end{enumerate}
}

%% memo:
\memoanswer{
本题采用半偏法测电表电阻，实验时为了测量更加准确，要求回路中的总电流尽可能保持不变，所以应使用电阻值比较大的滑动变阻器，这样可以尽可能使回路中的总电流保持不变；同时为了使电流表能够达到满偏，还需要选择合适的电源，若选择 $4\UV$ 的电源和 $10\UkO$ 的滑动变阻器，则电路中的最小电流为 $I = \frac{4\UV}{10\UkO} = 400\UmuA$，此时调节滑动变阻器，能够使电流表满偏（如果选择 A 则无法使电流表达到满偏），而且对电路中的电流控制比较好，所以电源选择 B，滑动变阻器选择阻值范围为 $10\UkO$ 的 D。
}

\item
%% number: 32
%% paperName: 1991年湖南云南海南高考第 32 题 7 分
%% typeId: 6
%% type: 计算题
%% chapter: 磁场
%% point: 带电粒子在磁场中的运动
%% method: 无
%% score: 7
%% degree: 650
%% duplicateId: 0
%% body:
一个初速度为零的带电粒子，电量为 $q$，质量为 $m$，经电势差为 $U$ 的电场区域加速后，射入磁感应强度为 $B$ 的匀强磁场中，粒子速度方向与磁场方向垂直。求粒子在磁场中运动轨道的半径。
%% answer: 
\jdanswer{
$\frac{1}{B}\sqrt{\frac{2mU}{q}}$
}

%% memo:
\memoanswer{
带电粒子在电场中运动，由动能定理得

$qU = \frac{1}{2}mv^{2}$ ①，

带电粒子在磁场中受洛伦兹力，做匀速圆周运动，由牛顿第二定律得

$qvB = m\frac{v^{2}}{r}$ ②，

由 ①② 两式解得

$r = \frac{1}{B}\sqrt{\frac{2mU}{q}}$ ③。
}

\item
%% number: 33
%% paperName: 1991年湖南云南海南高考第 33 题 7 分
%% typeId: 6
%% type: 计算题
%% chapter: 电磁感应
%% point: 远距离输电
%% method: 无
%% score: 7
%% degree: 550
%% duplicateId: 0
%% body:
有一理想的单相变压器，原、副线圈的匝数比为 100，原线圈上所加电压为 $23\UkV$，副线圈通过总电阻为 $2\UO$ 的供电导线向用户供电，用户用电器得到的电压是 $220\UV$，求供电导线上损耗的功率。
%% answer: 
\jdanswer{
$50\UW$
}

%% memo:
\memoanswer{
设变压器副线圈两端的电压为 $U_{2}$，输电线路损耗电压为 $\Delta U$，电器得到的电压为 $U_{3}$，供电导线上损耗功率为 $P$，由题意可得，

$\frac{U_{1}}{U_{2}}=\frac{23\times10^{3}\UV}{U_{2}}=100$，解得 $U_{2}=230\UV$，

则 $\Delta U = U_{2} - U_{3} = 230 - 220 = 10\UV$，

故 $P=\frac{\Delta U^{2}}{R}=\frac{10^{2}}{2}=50\UW$。
}

\item
%% number: 34
%% paperName: 1991年湖南云南海南高考第 34 题 8 分
%% typeId: 6
%% type: 计算题
%% chapter: 运动和力
%% point: 动量和能量
%% method: 无
%% score: 8
%% degree: 550
%% duplicateId: 0
%% body:
一静止的硼核 \ce{^{10}_{5}B} 吸收一个慢中子（速度可忽略）后，转变成锂核（\ce{^{7}_{3}Li}）并发射出一个粒子。已知该粒子的动能为 $1.8\UMeV$，求锂核的动能。
%% answer: 
\jdanswer{
$1.03\UMeV$
}

%% memo:
\memoanswer{
由核反应方程可知发射出粒子为 \ce{^{4}_{2}He} 核，以 $m_{1}$、$v_{1}$ 和 $E_{k1}$ 表示 \ce{^{4}_{2}He} 核的质量、速度和动能，$m_{2}$、$v_{2}$、$E_{k2}$ 表示 \ce{^{7}_{3}Li} 核的质量、速度和动能，由动量守恒得 $m_{1}v_{1}=m_{2}v_{2}$ ①，

又有 $E_{k1}=\frac{1}{2}m_{1}v_{1}^{2},\quad E_{k2}=\frac{1}{2}m_{2}v_{2}^{2}$ ②，

由 ② 和 ① 解得 $E_{k2}=\frac{m_{1}}{m_{2}}E_{k1}$ ③，

\ce{^{4}_{2}He}、\ce{^{7}_{3}Li} 的质量比 $\frac{m_{1}}{m_{2}}$ 等于它们的质量数之比为

$\frac{m_{1}}{m_{2}}=\frac{4}{7}$ ④，

代入 ③ 解得 $E_{k2}=\frac{4}{7}\times1.8=1.03\UMeV$。
}

\item
%% number: 35
%% paperName: 1991年湖南云南海南高考第 35 题 12 分
%% typeId: 6
%% type: 计算题
%% chapter: 气体性质
%% point: 气体多过程
%% method: 临界
%% score: 12
%% degree: 350
%% duplicateId: 0
%% body:
如图所示，一直立的气缸，由截面积不同的两圆筒联接而成。活塞 A、B 用一长为 $2l$ 的不可伸长的细线连接，它们可在筒内无摩擦地上下滑动。A、B 的截面积分别为 $S_{A} = 20\Ucmq$，$S_{B} = 10\Ucmq$。A、B 之间有一定质量的理想气体。A 的上方和 B 的下方都是大气，大气压强始终保持为 $1.0\times10^{5}\UPa$。
\begin{enumerate}
	\item
	当气缸内气体的温度为 $600\UK$、压强为 $1.2\times10^{5}\UPa$ 时，活塞 A、B 的平衡位置如图所示。已知活塞 B 的质量 $m_{B} = 1\Ukg$，求活塞 A 的质量 $m_{A}$。（计算时重力加速度取 $g = 10\Umsq$）
	\item
	已知当气缸内气体温度由 $600\UK$ 缓慢降低时，活塞 A 和 B 之间的距离保持不变，并一起向下缓慢移动（可认为两活塞仍处在平衡状态），直到活塞 A 移到两圆筒的联接处。若此后气体温度继续下降，直到活塞 A 和 B 之间的距离开始小于 $2l$ 为止。试分析在降温的整个过程中，气缸内气体压强的变化情况，并求出气体的最低温度。
\end{enumerate}
\onepicture[align=right, w=4.5cm]{figs/35.png}
%% answer: （1）1 kg；（2）300 K
\jdanswer{
\begin{enumerate}
	\item
	$1\Ukg$

	\item
	$300\UK$
\end{enumerate}
}

%% memo:
\memoanswer{
（1）对活塞 A 受力分析，A 受四个力的作用，由平衡条件得

$p_{0}S_{A} + m_{A}g + F_{1} - p_{1}S_{A} = 0$ ①，

对活塞 B 受力分析，B 受四个力的作用，由平衡条件得

$p_{1}S_{B}+m_{B}g-F_{1}-p_{0}S_{B}=0$ ②，

其中 $F_{1}$ 为细线的拉力，$p_{1}$ 为气缸中气体的压强。

$p_{0}=1.0\times10^{5}\UPa,\ p_{1}=1.2\times10^{5}\UPa$，

将 ①② 式联立解得

$m_{A} + m_{B} = 2\Ukg$，即 $m_{A} = 1\Ukg$。

（2）由题意，已知当气体温度缓慢降低时，A、B 两活塞间的距离保持不变，且处在平衡状态。在活塞 A 移到两圆筒的联接处之前，任何时刻作用于两活塞组成的系统的合力为零，即

$p_{0}S_{A}-p_{0}S_{B}+m_{A}g+m_{B}g-p_{2}S_{A}+p_{2}S_{B}=0$ ③，

式中 $p_{2}$ 为气体的压强。

由此可知 $p_{2}=p_{1}$，

即在温度由 $600\UK$ 下降到活塞 A 到达联接处的过程中，气缸内气体压强保持不变，为等压降温过程。活塞 A 到达联接处时，气缸内气体的温度 $T_{2}$，由盖-吕萨克定律得

$\frac{S_{A}l+S_{B}l}{T_{1}}=\frac{2S_{B}l}{T_{2}}$ ④，

在温度为 $T_{2}$ 时，活塞 B 处于静止状态，作用于活塞 B 的合力为零，即

$p_{0}S_{B}-pS_{B}+F-m_{B}g=0$ ⑤，

式中 $F$ 为线对活塞 B 的拉力，$p$ 为气缸内气体的压强。

当温度由 $T_{2}$ 继续缓慢下降时，$p$ 变小，$F$ 亦变小，但活塞 B 仍静止直到 $F$ 减小至零为止，在这过程中气缸内气体经历等容降温过程。

当 $F=0$ 后，若继续降温，活塞 B 将向上移动，两活塞间距离开始小于 $2l$ 的温度就是 $F=0$ 时气体的温度 $T_{3}$，此时气缸内气体的压强 $p_{3}$，可由 ⑤ 式求得

$p_{3}=p_{0}-\frac{m_{B}g}{S_{B}}$ ⑥，

由查理定律得

$\frac{p_{2}}{p_{3}}=\frac{T_{2}}{T_{3}}$ ⑦，

由 ④⑥⑦ 式解得 $T_{3}=300\UK$。
}

\end{enumerate}

\end{document}
